% OWNER: Arizona. Target: about 0.65 page; select findings, do not copy all extras.
\subsubsection{Accent association}

\begin{table}[t]
\caption{Accent, input-modality and source-wording diagnostics for Direct. All checks are
post hoc and leave the archived predictions unchanged. Intervals are 95\%:
sentence-clustered regression (word frequency) and Wilson (special ID, word
removal).}
\label{tab:further-diagnostics}
\centering
\footnotesize
\begin{tabularx}{\columnwidth}{@{}>{\raggedright\arraybackslash}p{1.25cm}
  >{\raggedright\arraybackslash}p{2.05cm}>{\raggedright\arraybackslash}X@{}}
\toprule
Check & Input; sample & Finding \\
\midrule
Accent & Archived output; 4,200 clips &
  Joint accent test $\chi^2(6)=5.34$, $p=0.50$ \\
\addlinespace[2pt]
Text input & Reference transcript; 4,200 clips &
  449 of 508 speech-marked clips marked again (88.4\%) \\
\addlinespace[2pt]
Special ID & Speech, token IDs kept; 700 clips &
  ID~1 in 74 clips (10.6\% [8.5, 13.1]); none has the marker \\
\addlinespace[2pt]
Word frequency & Archived output; 4,200 clips &
  OR 0.52 [0.48, 0.57] per Zipf unit of the rarest content lemma,
  length-adjusted \\
\addlinespace[2pt]
Word removal & Reference transcript; 385 sentences &
  Re-translation without the marker in 310 (80.5\% [76.3, 84.2]) after
  selected-word vs.\ 38 (9.9\% [7.3, 13.3]) after random-word deletion \\
\bottomrule
\end{tabularx}
\end{table}

The group marker rates reported above span about three percentage
points, and the speaker-clustered joint test did not detect an association
between accent and marker presence (Table~\ref{tab:further-diagnostics});
sentence clustering or length adjustment does not change this ($p=0.45$ and
$0.48$, respectively). Of the 370 normalized sentences recorded more than
once, six were marked in some recordings but not others.

\subsubsection{Input modality}
\label{sec:further-diagnostic-results}
With reference-transcript input, the marker recurs in most speech-marked
clips and appears in only 79 of the 3,692 speech-unmarked clips (2.1\%), so
the anomaly is not exclusive to speech input. Whisper-transcript input gives
a similar overall rate (11.9\%, versus 12.6\% with reference transcripts).
Marker status differs between speech and reference-transcript input for 138
clips (59 marked only with speech input, 79 only with reference-transcript
input). In the
700-clip speech re-run, 694 outputs reproduced their archived strings exactly,
and none of the 74 clips with special UNK ID~1 showed the visible marker in
either run. Outputs with hidden ID~1 can therefore remain in the marker-free
subset of Section~\ref{sec:synthesis-results}.

\subsubsection{Source wording}
Marker presence is associated with source vocabulary
that is rare in general English. The rate falls from 31.3\% (82 of 262 clips)
when the rarest content lemma has Zipf $\le 2.5$ to 5.1\% (119 of 2,323) when
it exceeds 4 (including imputed clips). After length adjustment, each
tenfold increase in that lemma's frequency is associated with about half the
odds of a marker (Table~\ref{tab:further-diagnostics}). Proper nouns show no
clear association (OR 0.95 [0.75, 1.21], adjusted for length and rarest-lemma
frequency). Among the 385 sentences with marked text-input translations,
re-translation without the selected word was marker-free for 80.5\% of them,
against 9.9\% without a random other word
(Table~\ref{tab:further-diagnostics}). The gap remains in the 161 sentences
whose random word is not a spaCy stop word (78.3\% versus 13.0\%). Selected words are rarer than the other non-stop words
(median surface-form Zipf 3.8 versus 4.6). Because each was chosen to
maximize the drop in summed marker-start probability at one decoding step,
the deletion comparison measures sensitivity to the targeted deletion and
does not isolate a causal effect of word rarity.

Marker occurrence is associated with rarer source wording, while no
association with accent was detected; its internal cause remains unresolved. The text-input and token-ID checks also show that
neither speech input nor generation of special UNK ID~1 is required for the
visible marker to occur. The word-removal experiment measures marker
presence rather than translation quality. Section~\ref{sec:synthesis-results} examines whether
the system ranking persists after visibly marked clips are excluded.
